\subsection{挠子的积分}\etiquette[1.3]{XVIII.subsec1.3}

\subsubsection{}\etiquette[1.3.1]{XVIII.subsubsec1.3.1} 设 $S$ 为概形，$G$ 为 $S$ 的大 fppf 景上的阿贝尔层，满足以下三个条件。

\paragraph{}\etiquette[1.3.1.1]{XVIII.par1.3.1.1} \begin{itemize} \item[(i)] \emph{对\pageoriginale每个满足 $r\geq 1$ 的射影空间 $p:\PP^r_{S}\rightarrow S$，在 $S$ 的大 fppf 景上}（\hyperrefext[XVII][XVII.subsec0.10]{0.10}）\emph{有} \begin{gather*}
G\xrightarrow{\sim} p_{\ast}p^{\ast}G\quad\mbox{且}\\
\Hhom(\GG_{m},G)\xrightarrow{\sim} R^{1}p_{\ast}p^{\ast}G
\end{gather*}

\item[(ii)] \emph{对每个 $n$ 及每个光滑仿射 $S$-态射 $f:X\rightarrow T$，若 $K_{1}$ 和 $K_{2}$ 为 $\Sym^{n}_{T}(X)$ 上的两个 $G$-挠子，则集合 $\Hom(K_{1},K_{2})$ 可与从 $K_{1}$ 在 $(X/T)^{n}$ 上的逆像到 $K_{2}$ 在 $(X/T)^{n}$ 上的逆像的对称态射所成的集合等同。}

此条件等价于以下两个条件同时成立：

\item[${\rm (ii')}$] $\Hom_{S}(\Sym^{n}_{T}(X),G)\xrightarrow{\sim}\Hom_{S}((X/T)^{n},
G)^{\sigma_n}$，以及

\item[${\rm (ii'')}$] 若 $\Sym^{n}_{T}(X)$ 上的挠子 $K$ 在 $(X/T)^{n}$ 上的逆像有对称截面，则 $K$ 平凡。

\item[(iii)] $G$ \emph{或者满足条件}（\hyperrefext[XVII][XVII.par6.3.21.1]{6.3.21.1}），\emph{或者同时满足条件}（\ref{XVII.par6.3.24.1}）\emph{及}（\ref{XVII.par6.3.24.2}）。 \end{itemize}

于是 \hyperref[XVIII.thm1.2.2]{1.2.2} 的结论成立（\cf \hyperref[XVIII.par1.2.2.1]{1.2.2.1}），\hyperrefext[XVII][XVII.subsubsec6.3.26]{6.3.26} 的公式也成立。

条件 (i) 已在 \hyperref[XVIII.thm1.2.2]{1.2.2} 及 \hyperref[XVIII.var1.2.15]{1.2.15} 中讨论。当 $G$ 可表示且为有限表现，或为 $S$ 的小平展景上的层的逆像，或为平坦拟凝聚层时，条件 (ii) 均成立。因此，以下各情形均满足条件 \hyperref[XVIII.par1.3.1.1]{1.3.1.1}。

\paragraph{}\etiquette[1.3.1.2]{XVIII.par1.3.1.2} \begin{enumerate} \renewcommand{\theenumi}{\alph{enumi}} \renewcommand{\labelenumi}{(\theenumi)} \item $G$ 光滑且为有限表现（\hyperref[XVIII.thm1.2.2]{1.2.2} 及 \hyperrefext[XVII][XVII.par6.3.21.1]{6.3.21.1}）。

\item $G$ 为 $S$ 的小平展景上的挠层的逆像（\hyperref[XVIII.var1.2.15]{1.2.15} (ii) 及 \hyperrefext[XVII][XVII.par6.3.21.1]{6.3.21.1}）。

\item $G$ 仿射，且为光滑群之间的满态射的核（\hyperref[XVIII.thm1.2.2]{1.2.2} 及 XVII \hyperref[XVII.eq6.3.3.1]{6.3.3.1}）。

\item $G$ 由 $S$ 上的平坦拟凝聚层定义（\hyperref[XVIII.var1.2.15]{1.2.15} (i) 及 \hyperrefext[XVII][XVII.par6.3.21.1]{6.3.21.1}）。 \end{enumerate}

\subsubsection{}\etiquette[1.3.2]{XVIII.subsubsec1.3.2} 设\pageoriginale $f:X\rightarrow S$ 为概形 $S$ 上的光滑射影曲线，纤维非空且几何连通。设 $\Ll$ 为 $X$ 上的可逆层，它在以下意义下“充分”相对丰沛。

\paragraph{}\etiquette[1.3.2.1]{XVIII.par1.3.2.1} 对 $S$ 的每个几何点 $s$，若 $X_{s}$ 的亏格为 $g$，则 \begin{gather*}
\deg_{X_{s}}(\Ll)>2g-2,\\
\deg_{X_{s}}(\Ll)\geq 1\quad\mbox{\rm 若}\quad g=0,\quaet
\deg_{X_{s}}(\Ll)\geq 2\quad\mbox{\rm 若}\quad g=1
\end{gather*}

此时有 $R^{1}f_{\ast}\Ll=0$，而 $f_{\ast}\Ll$ 局部自由，其构造与每个基变换可交换，且处处的秩 $\geq 2$。若 $f_{\ast}\Ll$ 的一个截面 $s$ 在任何点都不为零，将 $s$ 看作 $X$ 上的 $\Ll$ 的截面时，其零点概形 $D_{s}$ 是一个相对除子（\EGA IV 21.15.2）；由于 $X/S$ 的相对维数为 1，它在 $S$ 上有限局部自由。有典范同构 $$
s:\Oo(D_{s})\xrightarrow{\sim}\Ll{\quoi}.
$$ 除子 $D_{s}$ 对 $s$ 的依赖仅差乘以 $\Oo^{\ast}_{S}$ 的一个截面。因此，上述构造在由 $X$ 经基变换 $$
p:\PP(f_{\ast}(\Ll)^{\vee})\longrightarrow S{\quoi}.
$$ 所得的概形上给出一个相对除子 $Z_{\Ll}$。

设 $K$ 为 $G_{X}$ 下的挠子，$K|Z_{\Ll}$ 为它在 $Z_{\Ll}$ 上的逆像。此挠子从 $Z_{\Ll}$ 到 $\PP(f_{\ast}(\Ll)^{\vee})$ 的迹（\hyperrefext[XVII][XVII.subsubsec6.3.26]{6.3.26}），是 $\PP(f_{\ast}(\Ll)^{\vee})$ 上的一个挠子 $K_{\Ll}$： $$
K_{\Ll}=\Tr_{Z_{\Ll}/\PP}(K|Z_{\Ll}){\quoi}.
$$

由 \hyperref[XVIII.thm1.2.2]{1.2.2}，存在一个在唯一同构意义下唯一的三元组，由 $S$ 上的一个 $G$-挠子 $\langle \Ll$, $K$]、一个态射 $\varphi_{\Ll,K}:\GG_{m}\rightarrow G$ 及一个同构 $$
K_{\Ll}\simeq
p^{\ast}\langle \Ll,K]+\varphi_{\Ll,K}\Oo(1){\quoi}.
$$ 组成。

\subsubsection{}\etiquette[1.3.3]{XVIII.subsubsec1.3.3} 若\pageoriginale将挠子的合成用加法记号表示，则有典范同构及等式 \begin{align*}
& \langle \Ll,K_{1}+K_{2}] \simeq
\langle \Ll,K_{1}]+\langle \Ll,K_{2}]\tag{1.3.3.1}\\
& \varphi_{\Ll, K_{1}+K_{2}}=\varphi_{\Ll,K_{1}}+\varphi_{\Ll,K_{2}}{\quoi}.\tag{1.3.3.2}
\end{align*}\etiquette[1.3.3.1]{XVIII.eq1.3.3.1} \etiquette[1.3.3.2]{XVIII.eq1.3.3.2}

设 $\Ll_{1}$ 和 $\Ll_{2}$ 为 $X$ 上的两个满足（\hyperref[XVIII.par1.3.2.1]{1.3.2.1}）的可逆层，并设 $\Ll=\Ll_{1}\otimes \Ll_{2}$。截面的张量积定义一个态射 $$
\xymatrix@R=1.5cm@C=.5cm{
 & **[r] \PP
  (f_{\ast}(\Ll_1)^{\vee}) \ar[ld]^{\pr_1}& \times_S \ar[dd]^{p_{12}}& **[l]\PP
  (f_{\ast}(\Ll_2)^{\vee})\ar[rr]^{\pi}\ar[rd]^{\pr_2} &     &
 \PP(f_{\ast}(\Ll)^{\vee}) \ar@{-}[dd]^{p}\\
**[r]\PP(f_*(\Ll_1)^{\vee})\ar[rrd]_{p_1}&          &                &                 &  **[l]\PP(f_{\ast}(\Ll_2)^{\vee}) \ar[lld]^{p_2}  &\\
 &          &  S     \ar@{=}[rrr]        &                 &      & S.
}
$$ 有 $\pi^{\ast}Z_{\Ll}\simeq Z_{\Ll_{1}}+Z_{\Ll_{2}}$，从而得到一个

同构（\hyperrefext[XVII][XVII.const6.3.27.1]{6.3.27.1}） \begin{equation*}
\pi^{\ast}K_{\Ll}=\pr^{\ast}_{1}K_{\Ll_{1}}+\pr^{\ast}_{2}K_{\Ll_{2}}{\quoi}.\tag{1.3.3.3}
\end{equation*}\etiquette[1.3.3.3]{XVIII.eq1.3.3.3}

有典范同构 $$
\pi^{\ast}\Oo(1)\simeq \pr^{\ast}_{1}\Oo(1)\otimes
\pr^{\ast}_{2}\Oo(1){\quoi},
$$ 因而可将 \eqhyperref[XVIII.eq1.3.3.3]{1.3.3.3} 改写为 \begin{gather*}
\pi^{\ast}(p^{\ast}\langle \Ll,K]+\varphi_{\Ll,K}\Oo(1))=
p^{\ast}_{12}\langle \Ll,K]+\pr^{\ast}_{1}\varphi_{\Ll,K}\Oo(1)+\pr^{\ast}_{2}\varphi_{\Ll,K}\Oo(1) \\
= p^{\ast}_{12}\langle \Ll_{1},K]+p^{\ast}_{12}\langle \Ll_{2},K]+\pr^{\ast}_{1}\varphi_{\Ll_{1},K}\Oo(1)+\pr^{\ast}_{2}\varphi_{\Ll_{2},K}\Oo(1){\quoi}.
\end{gather*}

两次应用 \hyperref[XVIII.thm1.2.2]{1.2.2}，得到一个典范同构 \begin{gather*}
\langle \Ll_{1}\otimes \Ll_{2}, K]=\langle \Ll_{1},K]+\langle \Ll_{2},K]\tag{1.3.3.4}
\end{gather*}\etiquette[1.3.3.4]{XVIII.eq1.3.3.4} 及\pageoriginale一个恒等式 $$
\varphi_{\Ll_{1}\otimes\Ll_{2},K}=\varphi_{\Ll_{1},K}=\varphi_{\Ll_{2},K}{\quoi}.
$$

由此可见 $\varphi_{\Ll,K}$ 仅依赖于 $K$，于是置 \begin{equation*}
\deg(K)=\varphi_{\Ll,K}{\quoi}.\tag{1.3.3.5}
\end{equation*}\etiquette[1.3.3.5]{XVIII.eq1.3.3.5} 因此， \begin{align*}
& K_{\Ll}\simeq
p^{\ast}\langle \Ll,K]+\deg(K)\Oo(1){\quoi}.\tag{1.3.3.6}\\
& \deg(K_{1}+K_{2})=\deg(K_{1})+\deg(K_{2}){\quoi}.\tag{1.3.3.7}
\end{align*}\etiquette[1.3.3.6]{XVIII.eq1.3.3.6} \etiquette[1.3.3.7]{XVIII.eq1.3.3.7}

设 $\mathcal{U}$ 为一个可逆层，它局部为 $f_{\ast}\Ll$ 的直和因子，定义 $\PP((f_{\ast}\Ll)^{\vee})$ 的一个截面 $u$。设 $D_{u}$ 为 $\mathcal{U}$ 的局部可逆截面定义的相对除子。 $$
u^{\ast}\Oo(1)\simeq \mathcal{U}^{\vee}
$$ 因此，由 \eqhyperref[XVIII.eq1.3.3.6]{1.3.3.6} 得到 \begin{equation*}
u^{\ast}K_{\Ll}=\Tr_{D_u/S}(K)\simeq
\langle \Ll,K]-\deg(K)\mathcal{U}{\quoi}\nde{记号 $\Tr_{D_u/S}(K)$ 将在下文 \eqhyperref[XVIII.eq1.3.5.0]{1.3.5.0} 中引入。}.\tag{1.3.3.8}
\end{equation*}\etiquette[1.3.3.8]{XVIII.eq1.3.3.8}

以下相容性的验证留给读者：

\begin{enonce*}{命题 1.3.4}\etiquette[1.3.4]{XVIII.prop1.3.4} 同构 \eqhyperref[XVIII.eq1.3.3.1]{1.3.3.1} 及 \eqhyperref[XVIII.eq1.3.3.4]{1.3.3.4} 的以下图式交换： $$
\xymatrix@C=.1cm@R=.1cm{
\langle \Ll_{1}\otimes \Ll_{2}, K_{1}+K_{2}] & \simeq & \langle
\Ll_{1}\otimes \Ll_{2}, K_{1}]+\langle \Ll_{1} \otimes \Ll_{2}, K_{2}]\\
 & & \rotatebox{90}{$\simeq$} \\
\rotatebox{90}{$\simeq$} & & \langle \Ll_{1}, K_{1}]+\langle\Ll_{2}, K_{1}]+\langle
    \Ll_{1}, K_{2}]+\langle \Ll_{2}, K_{2}]\\
 & & \rotatebox{-50}{$\Arrowvert$}\\
\langle\Ll_{1}, K_{1}+K_{2}]+\langle\Ll_{2},K_{1}+K_{2}] & \simeq  &
    \langle\Ll_{1}, K_{1}]+\langle \Ll_{1},K_{2}]+\langle\Ll_{2},K_{1}]+\langle\Ll_{2},K_{2}]. 
}
$$

%$$
%\xymatrix{
%\langle \Ll_1\otimes \Ll_2, K_1 + K_2] \ar@{-}[dd]&\simeq& \langle \Ll_1\otimes \Ll_2, K_1] + \langle \Ll_! \otimes
%  \Ll_2, K_2]\ar@{-}[d]\\
%  & & \langle \Ll_1\otimes \Ll_2, K_1] + \langle \Ll_! \otimes
%  \Ll_2, K_2] \ar@{=}[d]\\
%\langle \Ll_1\otimes \Ll_2, K_1] + \langle \Ll_! \otimes
%  \Ll_2, K_2] & \simeq& \langle \Ll_1\otimes \Ll_2, K_1] + \langle \Ll_! \otimes
%  \Ll_2, K_2].
%}
%$$


\end{enonce*}

此外，同构 \eqhyperref[XVIII.eq1.3.3.1]{1.3.3.1} 及 \eqhyperref[XVIII.eq1.3.3.4]{1.3.3.4} 与结合及交换同构相容，精确表述留给读者。于是，当 $\Ll$ 不\pageoriginale必满足条件（\hyperref[XVIII.par1.3.2.1]{1.3.2.1}）时，可以用公式 \eqhyperref[XVIII.eq1.3.3.4]{1.3.3.4} 通过“线性”延拓 $\langle \Ll,K]$ 的定义。

\setcounter{subsubsection}{4} \subsubsection{}\etiquette[1.3.5]{XVIII.subsubsec1.3.5} 对 $X$ 上的每个有效相对除子 $D$，置 \begin{equation*}
\Tr_{D/S}(K)=\Tr_{D/S}(K|D){\quoi}.\tag{1.3.5.0}
\end{equation*}\etiquette[1.3.5.0]{XVIII.eq1.3.5.0}

若 $\Oo(D)$ 满足（\hyperref[XVIII.par1.3.2.1]{1.3.2.1}），则 $\Oo(D)$ 的截面 \hyperref[XVIII.sec1]{1} 是 $f_{\ast}\Oo(D)$ 的一个截面，作为这样的截面，它生成一个局部为 $f_{\ast}\Oo(D)$ 的直和因子的可逆层，并定义该层的一个平凡化。因此，同构 \eqhyperref[XVIII.eq1.3.3.8]{1.3.3.8} 给出一个同构 \begin{equation*}
\langle \Oo(D),K]\simeq \Tr_{D/S}(K){\quoi}.\tag{1.3.5.1}
\end{equation*}\etiquette[1.3.5.1]{XVIII.eq1.3.5.1}

立即验证，图式 $$\vcenter{
\xymatrix{
  \langle \Oo(D) , K_1 + K_2] \ar[r]^{\sim} \ar@{-}[d]_{\text{\rotatebox{90}{$\sim$}}}& \Tr_{D/S}(K_1 + K_2)\ar@{-}[d]^{\text{\rotatebox{90}{$\sim$}}}\\
\langle \Oo(D),K_1] + \langle \Oo(D) ,K_2] \ar[r]^{\sim}& \Tr_{D/S}(K_1) + \Tr_{D/S}(K_2)
}}\leqno{(1.3.5.2)}
$$\etiquette[1.3.5.2]{XVIII.eq1.3.5.2} 及 $$\vcenter{
\xymatrix{
\langle \Oo(D_1 + D_2),K] \ar[r]^{\sim}
\ar@{-}[d]^{\text{\rotatebox{90}{$\sim$}}}& \Tr_{(D_1 +
  D_{2})/S}(K)\ar@{-}[d]^{\text{\rotatebox{90}{$\sim$}}\quad (XVII\
  \hyperref[XVII.const6.3.27.1]{6.3.27.1})}\\ 
\langle \Oo(D_1),K] + \langle \Oo(D_2),K] \ar[r]^{\sim}& \Tr_{D_1 /S}(K) +
    \Tr_{D_{2} /S} (K)
}}\leqno{(1.3.5.3)}
$$\etiquette[1.3.5.3]{XVIII.eq1.3.5.3} 交换。由线性，可对每个相对除子定义同构 \eqhyperref[XVIII.eq1.3.5.1]{1.3.5.1}。

\subsubsection{}\etiquette[1.3.6]{XVIII.subsubsec1.3.6} 设 $D$ 和 $E$ 为两个有效相对 Cartier 除子，$f$ 为满足 $$
\text{\rm div\,}(f)=D-E{\quoi},
$$ 的有理函数，即 $\Oo(D)$ 与 $\Oo(E)$ 之间的一个同构。$\Oo(D)$ 与 $\Oo(E)$ 之间的同构 $f$ 通过 \eqhyperref[XVIII.eq1.3.5.1]{1.3.5.1} 定义一个同构 \begin{equation*}
\langle f,K]:\Tr_{D/S}(K)\xrightarrow{\sim} \Tr_{E/S}(K){\quoi}.\tag{1.3.6.1}
\end{equation*}\etiquette[1.3.6.1]{XVIII.eq1.3.6.1} 有\pageoriginale以下公式，其中第二个由 \eqhyperref[XVIII.eq1.3.5.2]{1.3.5.2} 得到： \begin{align*}
& \langle fg,K]=\langle f,K].\langle g,K]\\
& \langle f,K_{1}+K_{2}]=\langle f,K_{1}]+\langle f,K_{2}]{\quoi}.
\end{align*}

若 $\lambda$ 为 $\Oo^{\ast}_{S}$ 的一个截面，则对 $X$ 上的每个可逆层 $\Ll$，乘以 $\lambda$ 是 $\Ll$ 的一个自同构。通过结构转移，此自同构在 \eqhyperref[XVIII.eq1.3.3.8]{1.3.3.8} 的两边诱导自同构，且它们在同构 \eqhyperref[XVIII.eq1.3.3.8]{1.3.3.8} 下相互对应。于是 $$
\langle \lambda, K]=\deg (K)(\lambda){\quoi}\;;
$$ 因而若 $\deg(K)\neq 0$，则 $\Tr_{D/S}(K)$ 与 $\Tr_{E/S}(K)$ 之间的同构 $\langle f,K]$ 依赖于 $f$ 的选取。

设 $\lambda\in \Gamma(S,\Oo^{\ast}_{S})$、$g\in \Gamma(S,G)$，设 $\Ll$ 为 $X$ 上的可逆层，$K$ 为 $X$ 上的 $G$-挠子。若 $\langle \lambda,g]$ 为由 $\Ll$ 的自同构 $\lambda$ 及 $K$ 的自同构 $g$ 通过结构转移得到的 $\langle \Ll,K]$ 的自同构，则将 $\langle \lambda,g]$ 与 $G$ 的一个截面等同后，它由 $\langle \lambda,g]=\langle \lambda,K]+\langle \Ll,g]$ 给出；由于 $\langle \Oo(D),g]=\Tr_{D/S}(g)$，前一公式给出 \begin{equation*}
\langle \lambda,g]=\deg(K)(\lambda)+\deg(\Ll).g{\quoi}.\tag{1.3.6.2}
\end{equation*}\etiquette[1.3.6.2]{XVIII.eq1.3.6.2}

\subsubsection{}\etiquette[1.3.7]{XVIII.subsubsec1.3.7} 设 $f:X\rightarrow S$ 为 $S$ 上的光滑射影曲线，$G$ 为 $S$ 上满足 \hyperref[XVIII.par1.3.1.1]{1.3.1.1} 的群。若 $X\xrightarrow{f_{0}}T\xrightarrow{\pi}S$ 为 $f$ 的 Stein 分解，则 $\pi$ 为 $S$ 的平展覆盖，且层 $f_{\ast}G_{m}$ 为环面 $\prod_{T/S}\GG_{m}$。在 $S$ 上局部地（对于平展拓扑），$X$ 为 $S$ 上一族具有非空几何连通纤维的曲线 $(X_{i})_{i\in I}$ 的不交并。若 $K$ 为 $G$ 下的挠子，$\Ll$ 为 $X$ 上的可逆层，则置 $$
\langle \Ll,K]=\sum_{i}\langle \Ll|X_{i},K|X_{i}]
$$ 并定义 $K:f_{\ast}\GG_{m}=\GG^{I}_{m}\rightarrow G$ 的次数，其坐标为 $\deg(K|X_{i}):\GG_{m}\rightarrow G$。局部构造可整体化，从而给出\pageoriginale函子及态射 \begin{align*}
& \langle \Ll,K]_{X/S}=\Tr_{T/S}(\langle \Ll,K]_{X/T})\tag{1.3.7.1}\\
& \deg(K):f_{\ast}\GG_{m}=\prod_{T/S}\GG_{m}\longrightarrow G:\deg_{X/S}(K)=\prod_{T/S}\deg_{X/T}(K){\quoi}.\tag{1.3.7.2}
\end{align*}\etiquette[1.3.7.1]{XVIII.eq1.3.7.1} \etiquette[1.3.7.2]{XVIII.eq1.3.7.2}

对于 $\Ll$ 的由 $\Gamma(T,\Oo^{\ast}_{T})$ 的一个元素 $\lambda_{0}$ 定义的自同构 $\lambda$，有 \begin{equation*}
\langle \lambda,K]=\deg(K)(\lambda){\quoi}.\tag{1.3.7.3}
\end{equation*}\etiquette[1.3.7.3]{XVIII.eq1.3.7.3}

$K$ 的总次数定义为复合 $$
\deg \tot(K):\GG_{m}\longrightarrow
f_{\ast}\GG_{m}\xrightarrow{\deg(K)}G{\quoi}.
$$ 由前述结果得到：

\begin{enonce*}[remark]{公式汇集 1.3.8}\etiquette[1.3.8]{XVIII.for1.3.8} 对每个光滑射影曲线 $f:X\rightarrow S$，有一个函子 $\langle~,~]$，其构造与每个基变换相容，它将 $X$ 上的可逆层 $\Ll$ 及 $f^{\ast}G$ 下的挠子 $K$ 对应到 $S$ 上 $G$ 下的挠子 $\langle \Ll,K]$。此函子赋有以下附加结构，并满足以下条件： \end{enonce*}

\setcounter{subsubsection}{8} \paragraph{}\etiquette[1.3.8.1]{XVIII.par1.3.8.1} $\langle \Ll,K]$ 关于 $\Ll$ 和 $K$ 双可加；双可加同构与结合及交换同构和基变换相容，满足相容性（\hyperref[XVIII.prop1.3.4]{1.3.4}），并且与下述同构（\hyperref[XVIII.par1.3.8.2]{1.3.8.2}）相容（通过 \hyperrefext[XVII][XVII.par6.3.25.3]{6.3.25.3} 及 \hyperrefext[XVII][XVII.const6.3.27.1]{6.3.27.1}）。

\paragraph{}\etiquette[1.3.8.2]{XVIII.par1.3.8.2} 对 $X/S$ 上的每个有效相对除子 $E$，有典范同构 $$
\langle \Oo(E),K]\simeq \Tr_{E/S}(K|E){\quoi}.
$$

\paragraph{}\etiquette[1.3.8.3]{XVIII.par1.3.8.3} 对 $\Ll$ 的自同构 $\lambda\in f_{\ast}\GG_{m}$，次数 \eqhyperref[XVIII.eq1.3.7.2]{1.3.7.2} 满足 $$
\langle \lambda,K]=\deg(K)(\lambda){\quoi}.
$$ 对 $\lambda\in \GG_{m}(S)$，总次数满足 $$
\langle \lambda,K]=\deg\tot(K)(\lambda){\quoi}.
$$ 同样，\pageoriginale对 $g\in G(S)$ 有 $\langle \Ll,g]=\deg(\Ll).g$。

由此推出典范同构： \begin{equation*}
\langle f^{\ast}\Ll,K]\simeq \deg\tot(K)(\Ll){\quoi}.\tag{1.3.8.4}
\end{equation*}\etiquette[1.3.8.4]{XVIII.eq1.3.8.4}

由定义立即可知，对每个同态 $\varphi:G\rightarrow G'$，有一个与上述资料相容的典范同构： \begin{equation*}
\langle \Ll,\varphi(K)]\simeq \varphi(\langle \Ll,K]){\quoi}.\tag{1.3.8.5}
\end{equation*}\etiquette[1.3.8.5]{XVIII.eq1.3.8.5}

\setcounter{paragraph}{5} \paragraph{}\etiquette[1.3.8.6]{XVIII.par1.3.8.6} 设 $u:X\rightarrow Y$ 为 $S$ 上光滑射影曲线之间的平坦态射。则有与双可加性及基变换相容的同构： \begin{itemize} \item[--] 对 $Y$ 上的可逆层 $\Ll$ 及 $Y$ 上的 $G_{Y}$-挠子 $K$，有 $$
\langle \Ll, u^{\ast}K]\xrightarrow{\sim}\langle N_{X/Y}(\Ll),K]{\quoi}.
$$

\item[--] 对 $X$ 上的可逆层 $\Ll$ 及 $X$ 上的 $G_{X}$-挠子 $K$，有 $$
\langle u^{\ast}\Ll,K]\xrightarrow{\sim}
\langle \Ll,\Tr_{X/Y}K]{\quoi}.
$$ \end{itemize} 此外，这些同构对曲线态射的复合 $uv$ 满足显然的相容性；对 $\Ll=\Oo(E)$，它们与同构（\hyperrefext[XVII][XVII.const6.3.27.2]{6.3.27.2}）等同： \begin{align*}
\Tr_{E/S}(u^{\ast}K) &\xrightarrow{\sim} \Tr_{u_{\ast}E/S}(K)\\
\Tr_{u^{\ast}E/S}(K) &\xrightarrow{\sim} \Tr_{E/S}\Tr_{X/Y}K{\quoi}.
\end{align*} 还须证明 \hyperref[XVIII.par1.3.8.6]{1.3.8.6}。归结为 $X$ 和 $Y$ 具有非空几何连通纤维的情形，并注意到：对 $\Ll$ 的每个作为 $\Ll$ 的直接像的截面不为零的截面 $s$，上述显式公式给出一个所需类型的同构 $\Psi_{s}$，它关于 $s$ 齐次，因而当 $\Ll$ 满足（\hyperref[XVIII.par1.3.2.1]{1.3.2.1}）时与 $s$ 无关。由此即得一般情形。

\subsubsection{}\etiquette[1.3.9]{XVIII.subsubsec1.3.9} 在 \hyperref[XVIII.subsubsec1.3.7]{1.3.7} 的假设下，已在（\hyperref[XVIII.for1.3.8]{1.3.8}）中看到，$X$ 上的每个 $G$-挠子 $K$ 定义一个函子 $\langle \ast,K]$，以下记为 $F(\ast)$，它赋有\pageoriginale以下附加资料，并满足以下条件： \begin{enumerate} \renewcommand{\theenumi}{\roman{enumi}} \renewcommand{\labelenumi}{(\theenumi)} \item $F$ 为 $S$ 上的 fppf 叠的态射（\hyperrefext[XVII][XVII.subsec0.10]{0.10} 及 \textsc{Giraud} \cite{sga4.exp.XVIII.bib.1}） \begin{itemize} \item[--] 其源为如下叠：对一个 $S$-概形 $T$，其 $T$-对象为 $X_{T}$ 上的可逆层；

\item[--] 其目标为 $G$-挠子的叠。 \end{itemize}

\item 函子 $F$ 赋有可加同构 $F(\Ll_{1}\otimes \Ll_{2})\xrightarrow{\sim}F(\Ll_{1})+F(\Ll_{2})$，这些同构具有函子性，与结合及交换同构相容，并且与基变换相容。 \end{enumerate}

\begin{enonce*}{定理 1.3.10}\etiquette[1.3.10]{XVIII.thm1.3.10} 在 \hyperref[XVIII.subsubsec1.3.7]{1.3.7} 的假设下，函子 $\Phi:K\mapsto \langle \ast,K]$ 是 $X$ 上的 $G$-挠子范畴与（\hyperref[XVIII.subsubsec1.3.9]{1.3.9}）中所考虑的叠态射范畴之间的等价。 \end{enonce*}

容易归结为 $X$ 具有非空连通纤维的情形。

设 $F$ 满足 \hyperref[XVIII.subsubsec1.3.9]{1.3.9} (i) 及 (ii)。对每个 $S$-概形 $T$ 及 $X_{T}/T$ 的每个截面 $t$，也用 $t$ 表示相对除子 $t(T)$。则 $F(\Oo(t))$ 为 $T$ 上的 $G_{T}$-挠子。当 $t$ 为由 $X$ 到 $X$ 的恒等映射定义的 $X$ 的取值于 $X$ 的“泛点”时，用 $\Psi(F)$ 表示这样得到的 $X$ 上的挠子。对上述每个截面 $t$，有 \begin{equation*}
F(\Oo(t))\simeq t^{\ast}\Psi(F){\quoi}.\tag{1.3.10.1}
\end{equation*}\etiquette[1.3.10.1]{XVIII.eq1.3.10.1}

同构 \hyperref[XVIII.par1.3.8.2]{1.3.8.2} 给出函子的同构 $\Psi\circ\Phi \simeq \Id$。来构造一个同构 $\Phi\circ\Psi\simeq \Id$。

对 $X_{T}$ 在 $T$ 上的每个截面 $t$，同构 \eqhyperref[XVIII.eq1.3.10.1]{1.3.10.1} 是一个同构 $$
F(\Oo(t))\xrightarrow[\alpha]{\sim}\langle \Oo(t),\Psi(F)]{\quoi}.
$$

由可加性，$X_{T}$ 在 $T$ 上的每族截面 $t_{1},\ldots, t_{n}$ 定义一个同构 $$
F(\Oo(\sum t_{i}))\xrightarrow[\alpha]{\sim}\langle \Oo(\sum
t_{i}),\Psi(F)]{\quoi},
$$ 由 \hyperref[XVIII.subsubsec1.3.9]{1.3.9} (ii)，此同构不依赖于各 $t_{i}$ 的顺序。

取\pageoriginale $T=(X/S)^{n}$，并取 $t_{1}\ldots t_{n}$ 为截面的泛 $n^{\text{\rm 元组}}$。若 $T'=\Sym^{n}_{S}(X)$，则在 $X_{T'}$ 上有一个除子 $\sum t_{i}$，其逆像为 $X_{T}$ 上的除子 $\sum t_{i}$。在 $T$ 上，$\alpha$ 为一个对称同构 $$
\alpha:F(\Oo(\sum t_{i}))\xrightarrow{\sim} \langle \Oo(\sum
t_{i}),\Psi(F)]
$$，它联系的挠子均为 $T'$ 上挠子的逆像。由假设（\hyperref[XVIII.par1.3.1.1]{1.3.1.1}）(ii)，此同构来自 $T'$ 上唯一的同构： $$
\alpha:F(\Oo(\sum t_{i}))\xrightarrow{\sim} \langle \Oo (\sum
t_{i}),\Psi (F)]\qquad (\text{\rm 在 } T'){\quoi}.
$$

由 \hyperrefext[XVII][XVII.prop6.3.9]{6.3.9}，对 $X$ 上的每个相对除子 $D$，有一个同构 \begin{equation*}
\alpha_{D}:F(\Oo(D))\xrightarrow{\sim}
\langle \Oo(D),\Psi(F)]{\quoi},\tag{1.3.10.2}
\end{equation*}\etiquette[1.3.10.2]{XVIII.eq1.3.10.2} 它关于 $D$ 可加，且其构造与每个基变换相容。

设 $\Ll$ 为 $X$ 上的在 \hyperref[XVIII.par1.3.2.1]{1.3.2.1} 意义下充分丰沛的可逆层。$f_{\ast}\Ll$ 的每个处处非零截面 $s$，采用 \hyperref[XVIII.subsubsec1.3.2]{1.3.2} 的记号，定义一个同构 $$
s:\Oo(D_{s})\xrightarrow{\sim}\Ll{\quoi},
$$ 从而由 \eqhyperref[XVIII.eq1.3.10.2]{1.3.10.2} 定义一个同构 \begin{equation*}
\alpha_{s}:F(\Ll)\xrightarrow{\sim}\langle \Ll,\Psi(F)]{\quoi}.\tag{1.3.10.3}
\end{equation*}\etiquette[1.3.10.3]{XVIII.eq1.3.10.3}

证明此同构不依赖于 $s$。事实上： \begin{enumerate} \renewcommand{\theenumi}{\roman{enumi}} \renewcommand{\labelenumi}{(\theenumi)} \item $\alpha_{s}$ 的构造与每个基变换可交换。

\item 存在同态 $\varphi:\GG_{m}\rightarrow G$，使 $$
\alpha_{\lambda s}=\varphi(\lambda).\alpha_{s}
$$ （由 \eqhyperref[XVIII.eq1.3.10.3]{1.3.10.3} 两边的函子性）。 \end{enumerate}

必有 $\varphi=0$，否则由拉回 $F(\Ll)-\langle \Ll,\Psi(F)]$ 得到的 $\PP((f_{\ast}\Ll)^{\vee})$ 上的挠子将具有次数 $\neq 0$。因此，各 $\alpha_{s}$ 定义一个从 $\PP((f_{\ast}\Ll)^{\vee})$ 到 $G$ 的态射，而由假设\pageoriginale（\hyperref[XVIII.par1.3.1.1]{1.3.1.1}）(i)，此态射为常值。因此可置 $\alpha_{\Ll}=\alpha_{s}$。

各同构 $\alpha_{\Ll}$ 因此构成函子 $F$ 与 $\langle *,\Psi(F)]$ 在充分丰沛层上的限制之间的一个函子同构。此同构与可加同构相容，因而可通过线性延拓到所有可逆层，从而完成同构 $$
\alpha:\Phi\circ\Psi\simeq \Id{\quoi}.
$$ 的构造。

